\section{Einleitung}

\subsection{Das Problem, das nicht verschwindet}
Warum spüren Sie Schmerz? Das ist die einfachste Frage der Welt. Und zugleich die schwerste. Wir alle wissen, was Schmerz ist. Wir alle haben ihn gespürt. Aber warum spüren wir ihn? Wir wissen, dass Nozizeptoren Signale an den Thalamus senden, dass der Kortex sie verarbeitet und dass Endorphine freigesetzt werden. All das ist eine funktionale Beschreibung. Sie sagt uns, was das Gehirn tut, aber sie sagt uns nicht, warum Sie es spüren.

Denken Sie an einen Computer. Er kann Informationen verarbeiten, ein Bild erkennen, Schach spielen. Aber er fühlt nichts. Er hat kein subjektives Erleben. Nichts tut ihm weh. Warum ist das Gehirn anders? Das ist das schwierige Problem des Bewusstseins, formuliert von David Chalmers im Jahr 1995. Seither hat die Neurowissenschaft enorme Fortschritte gemacht, doch das schwierige Problem ist unangetastet geblieben. Niemandem ist es gelungen zu erklären, warum Materie zu Erleben wird. 1998 wettete Christof Koch mit Chalmers, dass ein neuronales Korrelat des Bewusstseins innerhalb von 25 Jahren gefunden würde. 2023 räumte er ein, dass die Wette verloren ist.

\subsection{Dunkle Materie --- die genauen Zahlen}
Das Universum hat eine Zusammensetzung, die wir ziemlich gut kennen. Die Planck-Mission (2018) hat uns die genauen Zahlen geliefert. Nur 4,9\% des Universums sind gewöhnliche Materie --- Atome, Sterne, Planeten, Menschen. Die Dunkle Materie macht 26,8\% aus, die Dunkle Energie 68,3\%. Praktisch sind 95,1\% des Universums unbekannt.

Stellen Sie sich das Universum als eine Truhe mit 100 Münzen vor. Nur 5 Münzen sind sichtbar. Die übrigen 95 sind unsichtbar. Davon sind 27 Dunkle Materie und 68 Dunkle Energie. Von der Dunklen Materie wissen wir, dass sie existiert (die gravitativen Belege sind solide), dass sie kein Licht aussendet und keines absorbiert und dass sie uns durchdringt, ohne wechselzuwirken. Etwa 100 Millionen Teilchen durchqueren Ihren Körper in jeder Sekunde. Wir wissen nicht, woraus sie besteht.

Eine neuere Einschränkung (2025) von Bonvin et al. hat gezeigt, dass Dunkle Materie genau wie gewöhnliche Materie in Gravitationspotentialtöpfe fällt. Wenn es eine „fünfte Kraft`` gibt, so kann sie 7\% der Gravitationskraft nicht überschreiten.

\subsection{Was ist ART-12D?}
ART-12D geht von einer einfachen Idee aus: Das Universum hat mehr Dimensionen als die vier, die wir wahrnehmen. Diese zusätzlichen Dimensionen sind kompakt --- so klein, dass wir sie nicht bemerken. Die Dunkle Materie existiert in allen zwölf Dimensionen. Wenn sie in die 4D projiziert wird, erscheint sie als ein Feld, das das Universum erfüllt. Dieses Feld ist das, was das Gehirn als Bewusstsein empfängt.

Der Name stammt von: \textbf{Akashic} (das universelle Informationsfeld), \textbf{Resonance} (weil die Wechselwirkung eine Resonanz ist), \textbf{Theory} (ein theoretischer Rahmen) und \textbf{12D} (weil ich in zwölf Dimensionen arbeite).

Die drei Postulate lauten: 1. Die Dunkle Materie ist ein fundamentales Quantenfeld. 2. Das Gehirn ist ein geometrischer Empfänger. 3. Die Kopplung erfolgt über eine Yukawa-Wechselwirkung.

Denken Sie an ein Radio. Es erzeugt die Musik nicht, es fängt sie aus der Luft auf. Die Radiowellen existieren, ob das Radio eingeschaltet ist oder nicht. Wenn das Radio kaputtgeht, verschwindet die Musik nicht aus dem Universum --- nur ihr lokaler Empfang endet. Ebenso funktioniert das Gehirn genau wie ein Radio. Das Bewusstsein wird nicht von den Neuronen erzeugt, sondern von ihnen empfangen. Und was sie empfangen, ist das Feld der Dunklen Materie.

\subsection{Warum zwölf Dimensionen?}
Warum zwölf? Warum nicht 10 oder 11 oder 26? Die Antwort kommt aus der Stringtheorie. Stellen Sie sich eine Gitarrensaite vor. Sie kann auf mehrere Arten schwingen. Jede Art entspricht einem Ton. Ebenso kann ein fundamentaler String in mehreren Dimensionen schwingen. Die Zahl der Dimensionen bestimmt, welche Schwingungsarten möglich sind.

Die Geschichte der Dimensionen: Kaluza--Klein (5 Dimensionen, 1919--1926), Superstrings (10, 1970--1980), M-Theorie (11, Witten, 1995), F-Theorie (12, Vafa, 1996), ART-12D (12, Ofiteru, 2026). Zwölf ist die größte Zahl von Dimensionen, in der die Stringtheorie mathematisch konsistent ist. Mehr Dimensionen würden zu Inkonsistenzen führen, weniger zu einer unvollständigen Theorie.

\subsection{Wie sehen die kompakten Dimensionen aus?}
Stellen Sie sich einen Gartenschlauch vor. Aus der Ferne sieht er eindimensional aus --- er hat nur Länge. Aus der Nähe hat er Dicke --- er ist zweidimensional (die Oberfläche des Zylinders). Ebenso sieht unser Raum dreidimensional aus. Aus der Nähe jedoch (auf der Planck-Skala, $10^{-35}$ m) könnte er elf Dimensionen haben. Die acht kompakten Dimensionen sind in jeden Punkt des Raumes eingerollt.

Was bedeutet „eingerollt``? Es sind keine Orte, an die man reist. Es sind keine Parallelwelten, die man durch Reisen erreicht. Es sind lokale Eigenschaften jedes Punktes. Jeder Punkt der Raumzeit hat zusätzlich acht kleine Dimensionen, die an ihn „angehängt`` sind.

\subsection{Wie hängen die Dimensionen mit dem Bewusstsein zusammen?}
Die Dunkle Materie existiert in zwölf Dimensionen. Sie ist ein Skalarfeld $A(x,y)$, wobei $x$ die 4D-Koordinaten und $y$ die acht kompakten Koordinaten sind. Wenn wir die acht Dimensionen kompaktifizieren, zerlegt sich das Feld in Kaluza--Klein-Moden:
\begin{equation}
A(x,y) = \sum_n A_n(x) Y_n(y)
\end{equation}
Jede Mode hat eine Masse und eine charakteristische Frequenz:
\begin{equation}
m_n^2 = \frac{n^2}{R^2}
\end{equation}
Das Gehirn kann durch seine geometrische Struktur (die Mikrotubuli in den Neuronen) mit bestimmten Moden resonieren. Wenn es resoniert, wird das Feld in subjektives Erleben umgesetzt.

Warum zählt die Architektur und nicht die Größe? Das Kleinhirn hat 69 Milliarden Neuronen (80\% der Gesamtzahl), erzeugt aber kein Bewusstsein. Es hat eine \textit{Feedforward}-Architektur --- die Information fließt in eine einzige Richtung, wie bei einer Drahtantenne. Der Kortex hat nur 16 Milliarden Neuronen, ist aber der Sitz des Bewusstseins. Er hat eine \textit{rekurrente} Architektur --- die Information kreist in Schleifen, wie bei einer gut ausgerichteten Parabolantenne. Die Schlussfolgerung: Es zählt nicht die Zahl der Neuronen, sondern die Geometrie des Netzwerks.

\section{Methodik}
\subsection{Wie ich gearbeitet habe}
Ich bin kein professioneller Physiker. Ich habe kein Labor. Ich habe keine Förderung. Was ich habe, ist eine Bibliothek, Zeit und eine Frage, die mich nicht loslässt. Ich habe in vier Schritten gearbeitet. Der erste Schritt war die Formulierung der Hypothese, ausgehend von der Intuition, dass Bewusstsein empfangen und nicht erzeugt wird. Der zweite Schritt war die Formalisierung mit den Werkzeugen der bekannten Mathematik: Kaluza--Klein für zusätzliche Dimensionen, F-Theorie für zwölf Dimensionen, Feldtheorie für Lagrangedichten, IIT für integrierte Information. Der dritte Schritt war die Ableitung von Vorhersagen, also dessen, was wir beobachten müssten, wenn die Theorie richtig ist. Ich habe die heute testbaren Vorhersagen ausgewählt. Der vierte Schritt war der Vorschlag von Experimenten mit bescheidenen Mitteln, beginnend mit einem für 50.000 Dollar.

\subsection{Was mich geleitet hat}
Ich habe mich an vier Kriterien gehalten. Innere Konsistenz --- die Gleichungen dürfen sich nicht widersprechen. Falsifizierbarkeit --- die Theorie muss testbare Vorhersagen machen. Konvergenz --- die Theorie darf der etablierten Physik nicht widersprechen. Ehrlichkeit --- ich muss einräumen, dass sie spekulativ ist.

\subsection{Was ich verwendet habe}
Ich habe Kaluza--Klein (1919) für die Idee kompakter Dimensionen verwendet, die F-Theorie (Vafa, 1996) für den zwölfdimensionalen Rahmen, die IIT (Tononi, 2004) für die Messung integrierter Information, Orch OR (Penrose, 1996) für den Quantenkollaps in Mikrotubuli und die Quantenfeldtheorie für Lagrangedichten und Bewegungsgleichungen.

\subsection{Meine Grenzen}
Ich muss ehrlich über die Grenzen sein. Ich habe keine experimentellen Belege; die Theorie wurde nicht getestet. Ich habe kein Peer-Review; ich habe nicht in wissenschaftlichen Zeitschriften veröffentlicht. Ich habe keine akademische Anbindung; ich bin unabhängiger Forscher. Ich habe keine Förderung; meine Mittel sind begrenzt.
